\reading{MR-12}{Arbitrage Pricing with Term Structure Models}{defer}{3}
% ==========================================================

\begin{mrkeybox}[title={Learning objectives in this reading}]
\small
\begin{enumerate}[label=\textbf{\alph*.},leftmargin=1.7em]
\item Calculate the expected discounted value of a zero-coupon security using a binomial tree.
\item Construct and apply an arbitrage argument to price a call option on a zero-coupon security
      using replicating portfolios.
\item Define risk-neutral pricing and apply it to option pricing.
\item Explain the difference between true and risk-neutral probabilities and apply this
      difference to interest rate drift.
\item Explain how the principles of arbitrage pricing of derivatives on fixed-income securities
      can be extended over multiple periods.
\item Define option-adjusted spread (OAS) and apply it to security pricing.
\item Describe the rationale behind the use of recombining trees in option pricing.
\item Calculate the value of a constant-maturity Treasury swap, given an interest rate tree and
      the risk-neutral probabilities.
\item Evaluate the advantages and disadvantages of reducing the size of the time steps on the
      pricing of derivatives on fixed-income securities.
\item Evaluate the appropriateness of the Black-Scholes-Merton model when valuing derivatives on
      fixed-income securities.
\end{enumerate}
\end{mrkeybox}

% ---------------------------------------------------------
\lo{MR-12 a}{Calculate the expected discounted value of a zero-coupon security using a binomial tree.}

\begin{mrdefbox}[title={Binomial interest rate tree}]
A model where interest rates can move up or down --- two possibilities --- in each period,
creating a tree structure of possible interest rate paths over time.
\begin{enumerate}[label=\textbf{\alph*)}]
\item \term{Nodes:} points where interest rates can move up (U) or down (L).
\item \term{Forward rates:} rates at each node representing the next period's interest rate.
\item \term{No arbitrage:} the values for on-the-run issues generated using an interest rate
      tree should prohibit arbitrage opportunities.
\end{enumerate}
\end{mrdefbox}

\subsection{Valuing an option-free bond with the tree, using backward induction}
\begin{enumerate}[label=\textbf{\alph*)}]
\item \term{Backward induction:} start from the bond's maturity and work backwards to its
      current value.
\item At each node, the bond value is the \term{average of discounted future values} from the
      next two nodes.
\item The discount rate is the \term{forward rate} associated with the current node.
\end{enumerate}

\begin{center}
\begin{tikzpicture}[every node/.style={font=\scriptsize\sffamily},
  val/.style={rectangle, rounded corners=2pt, draw=FRMNavy!45, fill=PaleNavy,
    line width=0.6pt, inner sep=3.5pt, minimum width=1.25cm, align=center},
  leaf/.style={rectangle, rounded corners=2pt, draw=black!35, fill=FRMSoft,
    line width=0.6pt, inner sep=3.5pt, minimum width=1.1cm, align=center}]
  \node[val] (root) at (0,0) {925.21};
  \node[val] (u) at (3.3,1.5) {970.87};
  \node[val] (d) at (3.3,-1.5) {980.39};
  \node[val] (m) at (3.3,0) {975.61};
  \node[leaf] (l1) at (6.9,2.3) {1,000};
  \node[leaf] (l2) at (6.9,0.75) {1,000};
  \node[leaf] (l3) at (6.9,-0.75) {1,000};
  \node[leaf] (l4) at (6.9,-2.3) {1,000};
  \draw[FRMNavy!70, line width=0.8pt] (root) -- node[above left=-2pt, text=black]{$.8024$} (u);
  \draw[FRMNavy!70, line width=0.8pt] (root) -- node[below left=-2pt, text=black]{$.1976$} (d);
  \draw[black!45, line width=0.7pt] (u) -- node[above left=-2pt, text=black]{$q$} (l1);
  \draw[black!45, line width=0.7pt] (u) -- node[below left=-2pt, text=black]{$1-q$} (m);
  \draw[black!45, line width=0.7pt] (m) -- (l2);
  \draw[black!45, line width=0.7pt] (m) -- (l3);
  \draw[black!45, line width=0.7pt] (d) -- node[above left=-2pt, text=black]{$q$} (m);
  \draw[black!45, line width=0.7pt] (d) -- node[below left=-2pt, text=black]{$1-q$} (l4);
  \node[anchor=west, text=black!60, align=left, text width=4.3cm] at (8.3,1.2)
    {Do not forget to add a coupon or payoff into every node value derived.};
  \node[anchor=west, text=black!60, align=left, text width=4.3cm] at (8.3,-1.2)
    {If $q = 0.5$ it is a \textbf{real world} probability. If $q \neq 0.5$ it is
     \textbf{risk-neutral}: it equates the current price with the market price.};
\end{tikzpicture}
\end{center}
\figcap{Backward induction runs right to left. Each interior node is the probability-weighted
average of the two nodes to its right, discounted one period at that node's own forward rate.}

\begin{mrkeybox}[title={Two ways to build a risk-neutral interest rate tree}]
\begin{enumerate}[label=\textbf{\alph*)}]
\item \term{Method 1:} adjust the \textit{interest rates} on the tree to match market prices of
      on-the-run bonds, ensuring an arbitrage-free tree.
\item \term{Method 2:} adjust the \textit{probabilities} to match market prices of bonds derived
      from the model, resulting in risk-neutral probabilities.
\end{enumerate}
\end{mrkeybox}

% ---------------------------------------------------------
\lo{MR-12 b}{Construct and apply an arbitrage argument to price a call option on a zero-coupon security using replicating portfolios.}

\noindent\gapbadge\hspace{4pt}{\sffamily\fontsize{7.6}{9}\selectfont\color{black!55}%
The source notes state this objective but never construct a replicating portfolio. The
construction below is written from the GARP curriculum.}\par
\vspace{6pt}

\begin{mrgapbox}[title={Why this matters}]
Risk-neutral pricing is a shortcut. The \textit{reason} it works is the replicating portfolio
argument underneath it: if a portfolio of ordinary bonds reproduces the option's payoff in every
state, then by the law of one price the option must cost what that portfolio costs, or there is
a riskless arbitrage.
\end{mrgapbox}

\subsection{The setup}
Price a call option expiring in six months to purchase \$1,000 face value of a then six-month
zero at \$990.

\begin{itemize}
\item If on date 1 the six-month rate is \term{2.50\%}, a six-month zero sells for \$987.654, so
      the right to buy it at \$990 is \term{worthless}.
\item If the six-month rate is \term{1.50\%}, the price is \$992.556, so the option is worth
      $992.556 - 990 = \mathbf{\$2.556}$.
\end{itemize}

\subsection{Build a portfolio that reproduces those two payoffs}
Construct a portfolio on date 0 of six-month and one-year zeros such that it is worth \$0 in the
up state and \$2.556 in the down state. Let $F^5$ and $F^1$ be the face values of the six-month
and one-year zeros.

\begin{mrfmlbox}
\begin{align*}
F^5 + 0.987654\,F^1 &= \$0 && \text{(up state)}\\
F^5 + 0.992556\,F^1 &= \$2.556 && \text{(down state)}
\end{align*}
\mrvk{In the up state, the maturing six-month zero pays its face value $F^5$, and the once
one-year zero --- now a six-month zero --- is worth 0.987654 per dollar of face. The same
reading applies in the down state at the other price.}{}
\end{mrfmlbox}

Solving the two equations:
\[ F^5 = -\$515.00 \qquad\qquad F^1 = \$521.4375 \]

In words: the option is replicated by \term{buying} \$521.4375 face of one-year zeros and
\term{shorting} \$515.00 face of six-month zeros on date 0.

\subsection{Price it by the law of one price}
\[ 0.990099\,F^5 + 0.978842\,F^1 \;=\; -0.990099 \times 515.00 + 0.978842 \times 521.4375
   \;=\; \mathbf{\$0.504} \]

\begin{mrtrapbox}[title={What enforces the price}]
If the option cost \textit{less} than \$0.504, arbitrageurs could buy the option, short the
replicating portfolio, keep the difference, and have no future liabilities. If it cost
\textit{more}, they would short the option and buy the portfolio. Ruling out riskless arbitrage
is what pins the price at \$0.504.
\end{mrtrapbox}

% ---------------------------------------------------------
\lo{MR-12 c}{Define risk-neutral pricing and apply it to option pricing.}
\lo{MR-12 d}{Explain the difference between true and risk-neutral probabilities and apply this difference to interest rate drift.}

\begin{center}
\begin{tikzpicture}[every node/.style={font=\scriptsize\sffamily},
  val/.style={rectangle, rounded corners=2pt, draw=FRMNavy!50, fill=PaleNavy,
    line width=0.7pt, inner sep=4pt, minimum width=1.7cm, align=center},
  leaf/.style={rectangle, rounded corners=2pt, draw=black!35, fill=FRMSoft,
    line width=0.6pt, inner sep=4pt, minimum width=1.5cm, align=center}]
  \node[val] (root) at (0,0) {\$90.006\\ 4.5749\%};
  \node[val] (u) at (3.6,1.5) {\$93.299\\ 7.1826\%};
  \node[val] (d) at (3.6,-1.5) {\$94.948\\ 5.3210\%};
  \node[leaf] (l1) at (7.4,2.5) {\$100.00};
  \node[leaf] (l2) at (7.4,1.0) {\$100.00};
  \node[leaf] (l3) at (7.4,-1.0) {\$100.00};
  \node[leaf] (l4) at (7.4,-2.5) {\$100.00};
  \foreach \a/\b in {root/u, root/d, u/l1, u/l2, d/l3, d/l4}{
    \draw[FRMNavy!60, line width=0.8pt] (\a) -- node[midway, sloped, above, text=black]
      {$0.5$} (\b);}
  \node[anchor=north, text=black!60] at (0,-2.6) {\textbf{Today}};
  \node[anchor=north, text=black!60] at (3.6,-2.6) {\textbf{End of 1 year}};
  \node[anchor=north, text=black!60] at (7.4,-2.6) {\textbf{End of 2 years}};
  \node[anchor=west, text=FRMGold, align=left, text width=3.0cm] at (8.9,0)
    {$q = 0.5$, so these are \textbf{true} probabilities. It is a zero-coupon bond.};
\end{tikzpicture}
\end{center}

\begin{mrexbox}[title={Backward induction, three steps}]
\small
\begin{align*}
\textbf{Step 1:}\quad V_{1,U} &= \frac{(\$100 \times 0.5) + (\$100 \times 0.5)}{1.071826}
  = \$93.299\\[3pt]
\textbf{Step 2:}\quad V_{1,L} &= \frac{(\$100 \times 0.5) + (\$100 \times 0.5)}{1.053210}
  = \$94.948\\[3pt]
\textbf{Step 3:}\quad V_{0} &= \frac{(\$93.299 \times 0.5) + (\$94.948 \times 0.5)}{1.045749}
  = \$90.006
\end{align*}
\end{mrexbox}

\begin{mrkeybox}[title={True versus risk-neutral probabilities}]
\begin{itemize}
\item If $q = 0.5$, the probabilities are \term{true} (real world) probabilities.
\item If $q \neq 0.5$, the probabilities have been adjusted so that the model price equals the
      observed market price. Those adjusted probabilities are the \term{risk-neutral}
      probabilities, and pricing with them is \term{risk-neutral pricing}.
\item The same adjustment can be made on the rates instead of the probabilities, which is where
      the \term{interest rate drift} comes from: shifting the rates up or down across the tree
      is the alternative to reweighting the branches.
\end{itemize}
\end{mrkeybox}

% ---------------------------------------------------------
\lo{MR-12 e}{Explain how the principles of arbitrage pricing of derivatives on fixed-income securities can be extended over multiple periods.}

\begin{mrexbox}[title={Example --- European call on a 3-year, 7\% coupon bond, strike 100}]
The tree below is built in two passes. First fill in the bond price at every node by backward
induction; then compute the option value at every node from those bond prices.
\end{mrexbox}

\begin{center}
\begin{tikzpicture}[every node/.style={font=\tiny\sffamily},
  nd/.style={rectangle, rounded corners=2pt, draw=FRMNavy!45, fill=PaleNavy,
    line width=0.6pt, inner sep=3.5pt, text width=2.05cm, align=left}]
  \node[nd] (root) at (0,0) {Int.\ rate = 3.00\%\\ Bond = \$105.01\\ \textbf{Option = \$0.45}};
  \node[nd] (u) at (3.6,2.0) {Int.\ rate = 5.99\%\\ Bond = \$100.37\\ \textbf{Option = \$0.23}};
  \node[nd] (d) at (3.6,-2.0) {Int.\ rate = 4.44\%\\ Bond = \$103.65\\ \textbf{Option = \$1.20}};
  \node[nd] (uu) at (7.6,3.2) {Int.\ rate = 8.56\%\\ Bond = \$98.56\\ \textbf{Option = \$0.00}};
  \node[nd] (mm) at (7.6,0) {Int.\ rate = 6.34\%\\ Bond = \$100.62\\ \textbf{Option = \$0.62}};
  \node[nd] (dd) at (7.6,-3.2) {Int.\ rate = 4.70\%\\ Bond = \$102.20\\ \textbf{Option = \$2.20}};
  \draw[FRMNavy!65, line width=0.8pt] (root) -- node[sloped, above, text=black]{0.76} (u);
  \draw[FRMNavy!65, line width=0.8pt] (root) -- node[sloped, below, text=black]{0.24} (d);
  \draw[black!45, line width=0.7pt] (u) -- node[sloped, above, text=black]{0.60} (uu);
  \draw[black!45, line width=0.7pt] (u) -- node[sloped, above, text=black]{0.40} (mm);
  \draw[black!45, line width=0.7pt] (d) -- node[sloped, above, text=black]{0.60} (mm);
  \draw[black!45, line width=0.7pt] (d) -- node[sloped, below, text=black]{0.40} (dd);
  \node[anchor=north, text=black!60] at (0,-3.9) {\textbf{Today}};
  \node[anchor=north, text=black!60] at (3.6,-3.9) {\textbf{End of year 1}};
  \node[anchor=north, text=black!60] at (7.6,-3.9) {\textbf{End of year 2}};
\end{tikzpicture}
\end{center}
\figcap{A recombining tree: the middle node at year 2 is reached by an up-then-down move or a
down-then-up move, and carries one rate either way. Bond prices in plain type, option values in
bold.}

\subsection{Step 1 --- find the bond prices}
\begin{mrexbox}
\small
\textbf{Year 2 nodes} (one year to maturity, coupon 7, face 100):
\[ \text{Mid: } \frac{107}{1.0634} = \$100.62 \qquad\qquad
   \text{Bottom: } \frac{107}{1.0470} = \$102.20 \]

\textbf{Year 1 nodes} (add the \$7 coupon to each year-2 price before discounting):
\begin{align*}
\text{Top: } &\frac{(\$105.56 \times 0.6) + (\$107.62 \times 0.4)}{1.0599} = \$100.37\\[3pt]
\text{Bottom: } &\frac{(\$107.62 \times 0.6) + (\$109.20 \times 0.4)}{1.0444} = \$103.65
\end{align*}

\textbf{Today:}
\[ \frac{(\$107.37 \times 0.76) + (\$110.65 \times 0.24)}{1.03} = \$105.01 \]
\end{mrexbox}

\subsection{Step 2 --- find the option values}
\begin{mrexbox}
\small
\[ \frac{(\$0.00 \times 0.6) + (\$0.62 \times 0.4)}{1.0599} = \$0.23
   \qquad\longrightarrow\qquad
   \frac{(\$0.23 \times 0.76) + (\$1.20 \times 0.24)}{1.03} = \mathbf{\$0.45} \]
\end{mrexbox}

\begin{mrtrapbox}[title={The step people drop}]
The \$7 coupon must be added to each node's derived bond price \textit{before} it is discounted
back a period. \$105.56 is \$98.56 plus the coupon; \$107.62 is \$100.62 plus the coupon. The
option values, by contrast, carry no coupon.
\end{mrtrapbox}

% ---------------------------------------------------------
\lo{MR-12 f}{Define option-adjusted spread (OAS) and apply it to security pricing.}

\begin{mrdefbox}[title={Option-adjusted spread}]
The option-adjusted spread is the spread that makes the \term{model value} --- calculated as the
present value of projected cash flows --- equal to the \term{current market price}.
\end{mrdefbox}

\begin{mrexbox}[title={Applying OAS to the CMT swap below}]
The model price of the CMT swap was \$1,466.63. Now assume the market price was instead
\$1,464.40, which is \$2.23 \textit{less} than the model price. In this case, the OAS to be
added to each discounted risk-neutral rate in the CMT swap binomial tree turns out to be
\term{20 basis points}.
\end{mrexbox}

\begin{mrkeybox}[title={Direction}]
The market price sits \textit{below} the model price, so the discount rates must be
\textit{raised} to bring the model down to the market. A positive OAS is added to every rate in
the tree.
\end{mrkeybox}

% ---------------------------------------------------------
\lo{MR-12 g}{Describe the rationale behind the use of recombining trees in option pricing.}

\begin{mrdefbox}[title={Recombining and non-recombining}]
If the interest rate in the middle node is the \term{same} irrespective of an up-then-down or a
down-then-up move, the tree is \term{recombining}. Otherwise it is
\term{non-recombining}.
\end{mrdefbox}

\begin{center}
\begin{tikzpicture}[every node/.style={font=\scriptsize\sffamily},
  dot/.style={circle, fill=FRMNavy, inner sep=1.8pt},
  dotr/.style={circle, fill=FRMRed, inner sep=1.8pt}]
% recombining
\begin{scope}
  \node[dot] (a) at (0,0) {};
  \node[dot] (b) at (1.5,0.9) {}; \node[dot] (c) at (1.5,-0.9) {};
  \node[dot] (d) at (3.0,1.8) {}; \node[dot] (e) at (3.0,0) {}; \node[dot] (f) at (3.0,-1.8) {};
  \foreach \p/\q in {a/b, a/c, b/d, b/e, c/e, c/f}{
    \draw[FRMNavy!70, line width=0.8pt] (\p) -- (\q);}
  \node[anchor=west, text=FRMNavy] at (3.2,0) {one middle node};
  \node[anchor=north, text=black!70] at (1.8,-2.3) {\textbf{Recombining} --- 3 nodes at step 2};
\end{scope}
% non-recombining
\begin{scope}[shift={(8.2,0)}]
  \node[dotr] (a) at (0,0) {};
  \node[dotr] (b) at (1.5,0.9) {}; \node[dotr] (c) at (1.5,-0.9) {};
  \node[dotr] (d) at (3.0,1.8) {}; \node[dotr] (e1) at (3.0,0.45) {};
  \node[dotr] (e2) at (3.0,-0.45) {}; \node[dotr] (f) at (3.0,-1.8) {};
  \foreach \p/\q in {a/b, a/c, b/d, b/e1, c/e2, c/f}{
    \draw[FRMRed!70, line width=0.8pt] (\p) -- (\q);}
  \node[anchor=west, text=FRMRed] at (3.2,0) {two separate middle nodes};
  \node[anchor=north, text=black!70] at (1.8,-2.3)
    {\textbf{Non-recombining} --- 4 nodes at step 2};
\end{scope}
\end{tikzpicture}
\end{center}
\figcap{The rationale is arithmetic. A recombining tree grows to $n+1$ nodes after $n$ steps; a
non-recombining tree grows to $2^n$. Over 30 steps that is 31 nodes against more than a billion,
which is the difference between a tractable model and an impossible one.}

% ---------------------------------------------------------
\lo{MR-12 h}{Calculate the value of a constant-maturity Treasury swap, given an interest rate tree and the risk-neutral probabilities.}

A CMT swap is an agreement to swap a floating rate for a Treasury rate, such as the 10-year
rate. The rate tree can be used to price it.

\begin{mrfmlbox}
\[ \text{Payoff} \;=\; \left( \frac{\$1{,}000{,}000}{2} \right) \times
   \bigl( y_{CMT} - \text{fixed rate} \bigr) \]
\mrvk{$y_{CMT}$ = a \textit{semiannually compounded} yield of a predetermined maturity, observed
on the payment date;\quad the division by 2 is the semiannual accrual;\quad the fixed rate in
the worked example below is \textbf{7.00\%}.}{}
\end{mrfmlbox}

\begin{center}
\begin{tikzpicture}[every node/.style={font=\tiny\sffamily},
  nd/.style={rectangle, rounded corners=2pt, draw=FRMNavy!45, fill=PaleNavy,
    line width=0.6pt, inner sep=3.5pt, text width=1.85cm, align=left}]
  \node[nd] (root) at (0,0) {Rate = 7.00\%\\ \textbf{V = \$1,466.63}};
  \node[nd] (u) at (3.5,1.8) {Rate = 7.25\%\\ \textbf{V = \$2,697.53}};
  \node[nd] (d) at (3.5,-1.8) {Rate = 6.75\%\\ \textbf{V = $-$\$2,217.35}};
  \node[nd] (uu) at (7.2,3.0) {Rate = 7.50\%\\ payoff \$2,500};
  \node[nd] (mm) at (7.2,0) {Rate = 7.00\%\\ payoff \$0};
  \node[nd] (dd) at (7.2,-3.0) {Rate = 6.50\%\\ payoff $-$\$2,500};
  \draw[FRMNavy!65, line width=0.8pt] (root) -- node[sloped, above, text=black]{0.76} (u);
  \draw[FRMNavy!65, line width=0.8pt] (root) -- node[sloped, below, text=black]{0.24} (d);
  \draw[black!45, line width=0.7pt] (u) -- node[sloped, above, text=black]{0.60} (uu);
  \draw[black!45, line width=0.7pt] (u) -- node[sloped, above, text=black]{0.40} (mm);
  \draw[black!45, line width=0.7pt] (d) -- node[sloped, above, text=black]{0.60} (mm);
  \draw[black!45, line width=0.7pt] (d) -- node[sloped, below, text=black]{0.40} (dd);
  \node[anchor=north, text=black!60] at (0,-3.7) {\textbf{Today}};
  \node[anchor=north, text=black!60] at (3.5,-3.7) {\textbf{Six months}};
  \node[anchor=north, text=black!60] at (7.2,-3.7) {\textbf{One year}};
\end{tikzpicture}
\end{center}

\subsection{Step 1 --- calculate the payoffs, floating minus fixed}
\begin{mrexbox}
\small
\begin{align*}
\text{payoff}_{1,U} &= \tfrac{1{,}000{,}000}{2} \times (7.25\% - 7.00\%) = \$1{,}250\\
\text{payoff}_{1,L} &= \tfrac{1{,}000{,}000}{2} \times (6.75\% - 7.00\%) = -\$1{,}250\\
\text{payoff}_{2,U} &= \tfrac{1{,}000{,}000}{2} \times (7.50\% - 7.00\%) = \$2{,}500\\
\text{payoff}_{2,M} &= \tfrac{1{,}000{,}000}{2} \times (7.00\% - 7.00\%) = \$0\\
\text{payoff}_{2,L} &= \tfrac{1{,}000{,}000}{2} \times (6.50\% - 7.00\%) = -\$2{,}500
\end{align*}
\end{mrexbox}

\subsection{Step 2 --- value by backward induction}
\begin{mrexbox}
\small
\begin{align*}
V_{1,U} &= \frac{(\$2{,}500 \times 0.6) + (\$0 \times 0.4)}{1 + \tfrac{0.0725}{2}} + \$1{,}250
  = \$2{,}697.53\\[4pt]
V_{1,L} &= \frac{(\$0 \times 0.6) + (-\$2{,}500 \times 0.4)}{1 + \tfrac{0.0675}{2}} - \$1{,}250
  = -\$2{,}217.35\\[4pt]
V_{0} &= \frac{(\$2{,}697.53 \times 0.76) + (-\$2{,}217.35 \times 0.24)}{1 + \tfrac{0.07}{2}}
  = \mathbf{\$1{,}466.63}
\end{align*}
\end{mrexbox}

\begin{mrtrapbox}[title={Do not forget the mid values}]
Two things get dropped here. First, the \textbf{middle} year-1 payoff of \$0 still has to be
carried into both $V_{1,U}$ and $V_{1,L}$ --- it is not the same as leaving the branch out.
Second, the \textbf{six-month payoff itself} is added to the discounted value at that node:
$+\$1{,}250$ at the upper node, $-\$1{,}250$ at the lower one.
\end{mrtrapbox}

% ---------------------------------------------------------
\lo{MR-12 i}{Evaluate the advantages and disadvantages of reducing the size of the time steps on the pricing of derivatives on fixed-income securities.}

\begin{center}
\small
\begin{tabularx}{\linewidth}{@{}p{4.2cm} >{\raggedright\arraybackslash}X@{}}
\arrayrulecolor{FRMRule}\toprule
\rowcolor{FRMNavy} \hdr{Advantage} & \hdr{Disadvantage}\\
Smaller time steps \term{increase model accuracy}. Daily steps enhance accuracy.
  & They \term{add computational complexity} and require more computational resources.\\
\bottomrule
\end{tabularx}
\end{center}

% ---------------------------------------------------------
\lo{MR-12 j}{Evaluate the appropriateness of the Black-Scholes-Merton model when valuing derivatives on fixed-income securities.}

\begin{mrtrapbox}[title={Why BSM does not fit fixed income}]
The Black-Scholes-Merton model is \term{not suitable} for fixed-income securities due to
unrealistic assumptions:
\begin{enumerate}[label=\textbf{\alph*)}]
\item \term{No upper limit to the asset price.} Bonds have an upper limit.
\item \term{Constant risk-free rate.} Interest rates change, and bonds are affected by exactly
      that.
\item \term{Constant bond price volatility.} Bond price volatility is not constant --- it falls
      as the bond approaches maturity.
\end{enumerate}
\end{mrtrapbox}

\subsection{Bonds with embedded options}

\begin{center}
\small
\begin{tabularx}{\linewidth}{@{}p{2.4cm} >{\raggedright\arraybackslash}X@{}}
\arrayrulecolor{FRMRule}\toprule
\rowcolor{FRMNavy} \hdr{Callable bonds} & \hdr{}\\
a) & \term{Negative convexity} due to the callable feature.\\
\rowcolor{FRMSoft} b) & The price-yield curve \term{flattens} as yields decrease.\\
c) & Capital gains are \term{capped} by the call price.\\
\rowcolor{FRMSoft} d) & Increases \term{reinvestment risk}.\\
\midrule
\rowcolor{FRMNavy} \hdr{Putable bonds} & \hdr{}\\
a) & The putable feature affects the price-yield relationship.\\
\rowcolor{FRMSoft} b) & The put price serves as a \term{floor} value for the bond price.\\
c) & Putable bonds are \term{less sensitive} to yield increases compared to non-putable bonds at
  higher yields.\\
\bottomrule
\end{tabularx}
\end{center}

\begin{center}
\begin{tikzpicture}
\begin{axis}[
  width=11.0cm, height=6.3cm,
  xmin=0, xmax=12, ymin=55, ymax=155,
  xlabel={\scriptsize Yield (\%)}, ylabel={\scriptsize Bond price},
  xlabel style={font=\scriptsize}, ylabel style={font=\scriptsize},
  tick label style={font=\scriptsize},
  xtick={0,2,4,6,8,10,12}, ytick={60,80,100,120,140},
  legend style={font=\scriptsize\sffamily, draw=FRMRule, fill=white,
    at={(0.5,-0.30)}, anchor=north, legend columns=-1, column sep=10pt},
  axis line style={draw=black!55}, clip=false]
\addplot[black!45, line width=1.0pt, smooth, domain=1.35:12, samples=120]
  {60 + 420/(x+3.2)};
\addlegendentry{Option-free bond}
\addplot[FRMRed, line width=1.3pt, dashed, smooth, domain=1.35:12, samples=120]
  {min(115, 60 + 420/(x+3.2))};
\addlegendentry{Callable --- capped at the call price}
\addplot[FRMGreen, line width=1.3pt, dotted, smooth, domain=1.35:12, samples=120]
  {max(95, 60 + 420/(x+3.2))};
\addlegendentry{Putable --- floored at the put price}
\draw[FRMRed!55, line width=0.5pt, dashed] (axis cs:0,115) -- (axis cs:12,115);
\draw[FRMGreen!55, line width=0.5pt, dashed] (axis cs:0,95) -- (axis cs:12,95);
\node[anchor=west, font=\tiny\sffamily, text=FRMRed] at (axis cs:0.15,118) {call price};
\node[anchor=west, font=\tiny\sffamily, text=FRMGreen] at (axis cs:0.15,91) {put price};
\end{axis}
\end{tikzpicture}
\end{center}
\figcap{The embedded option shows up as a bend in the price-yield curve. The callable bond
cannot rise past the call price as yields fall, which is what negative convexity looks like. The
putable bond cannot fall past the put price as yields rise, which is why it is less sensitive to
yield increases at the high end.}


% ==========================================================
